\section{Related Work}
\label{sec:related}

\noindent\textbf{BPE as an optimisation problem.}
BPE was introduced for data compression \citep{gage1994} and adapted to subword
segmentation by \citet{sennrich2016}. \citet{zouhar2023} formalise it as the search for a
merge sequence that maximises the number of tokens saved, show that this objective is
submodular over merge sequences, and bound the approximation ratio of the greedy algorithm.
Unigram tokenization \citep{kudo2018} instead prunes a large vocabulary under a likelihood
objective. These formulations give tokenizer learning a clear objective, but it is a single
objective defined on the pooled corpus, with no notion of cost per language. In contrast, we
treat the per-language costs as a vector of objectives and study the family of
scalarisations that connects total compression to worst-case parity.

\noindent\textbf{Cross-lingual disparity and fair tokenizers.}
\citet{petrov2023} show that the tokenizers of deployed models encode the same content with
very different numbers of tokens across languages, and measure this cost on parallel text,
the unit we adopt. Parity-aware BPE \citep{foroutan2026} changes the merge rule: at every
step it selects the language with the worst compression on a multi-parallel development set
and merges the most frequent pair in that language. The same work proposes a hybrid
schedule and a moving-window variant to reach intermediate tokenizers.
\citet{kanjirangat2026} extend the parity rule to the $N$ worst-off languages and find that
increasing $N$ gives inconsistent effects across metrics. \citet{lee2026} compare equitable
tokenizers on Southeast Asian languages and find that Parity-aware BPE lies on the
efficiency--equity frontier but becomes less equitable as the vocabulary grows, since the
rule bounds only the worst case. A second line of work changes the input rather than the
rule, either by rebalancing the data on which the vocabulary is learned
\citep{chung2023unimax} or by replacing the byte encoding with a morphology-driven one
\citep{limisiewicz2024myte}. Nonetheless, the intermediate tokenizers in all these works are
obtained by schedules, by a choice of $N$, or by data weights, none of which is derived from
an objective over per-language costs. Our power-mean objective recovers both ends of this
spectrum as limits of one rule, and weights every language continuously by its current cost
instead of selecting a subset.

\noindent\textbf{Vocabulary allocation and Indian languages.}
Several works decide how much vocabulary each language should receive. VoCap
\citep{zheng2021vocap} allocates capacity per language from its average log probability,
and XLM-V \citep{liang2023xlmv} builds per-cluster vocabularies before merging them.
\citet{brahma2026} study 17 Indian languages in 11 scripts and set each cluster's share of
the vocabulary in proportion to its share of the union of monolingual vocabularies. On the
morphological axis, SuTRA \citep{rathore2026} penalises merges that cross morpheme
boundaries in Hindi, Marathi and Gujarati and releases gold segmentations for these
languages, while \citet{arnett2025} find across 70 languages that morphological alignment
explains little of the variance in downstream performance. We use the SuTRA languages as
our testbed, and treat allocation and morphology as complementary to the merge rule studied
here.

\noindent\textbf{Fairness through scalarisation.}
Power means are a standard way to scalarise competing objectives in multi-objective
optimisation \citep{miettinen1999}. In resource allocation, the same family appears as
$\alpha$-fair utility \citep{mo2000fair}, which moves from efficiency to max--min fairness
as $\alpha$ grows, and in economics as the Atkinson index of inequality
\citep{atkinson1970}. Our training algorithm also builds on the lazy greedy technique of
\citet{minoux1978} for submodular maximisation. To our knowledge, this family has not been
applied to the merge rule of a tokenizer, where it yields a closed-form weighting of pair
counts and an exact lazy trainer.